\begin{problem}[The structure sheaf of $\Spec\mathbb Z$]\label{prob:vi17-p07}
Compute $\OO(D(n))$, the stalk at $(p)$, the stalk at $(0)$, and the corresponding residue fields for $X=\Spec\mathbb Z$.
\end{problem}

\ifdefined\FullProblemDossiers
\paragraph{Why this problem matters.}
The arithmetic line shows in one example how localization, generic points, closed points, and varying residue fields interact.

\paragraph{Useful ideas.}
Use $\mathbb Z_n=\mathbb Z[1/n]$ and localize at prime ideals.

\paragraph{Guided hints.}
\begin{enumerate}
\item At $(p)$ invert all integers not divisible by $p$.
\item At $(0)$ invert every nonzero integer.
\item Then quotient the local rings by their maximal ideals.
\end{enumerate}

\begin{solution}
For every $n$, $\OO(D(n))=\mathbb Z[1/n]$.  At a closed point $(p)$,
\[\OO_{(p)}=\mathbb Z_{(p)},\qquad\kappa(p)=\mathbb F_p.\]
At the generic point,
\[\OO_{(0)}=\mathbb Q,\qquad\kappa(0)=\mathbb Q.\]
Thus a global integer specializes to its residue modulo $p$ at the closed point while retaining its rational value generically.
\end{solution}

\paragraph{What happened mathematically.}
One sheaf simultaneously records rational and modular arithmetic.

\paragraph{Extensions.}
Describe $D(n)$ as the arithmetic line with the prime divisors of $n$ removed.

\paragraph{Related problems.}
VI.17.P6 and VI.17.P8 develop neighboring aspects of the same localization-sheaf calculus.

\paragraph{Provenance.}
\textbf{New synthesis using recurring corpus examples for $\Spec\mathbb Z$.}
\else
\paragraph{Provenance.} \textbf{New synthesis using recurring corpus examples for $\Spec\mathbb Z$.}
\fi
